\chapter{DataFrames and data wrangling}
\label{ch:dataframes}

\begin{quote}
\emph{``Data wrangling is the process of transforming messy, real-world data into a clean, structured format suitable for analysis.''}
\end{quote}

% ============================================================
\section{The DataFrames.jl package}

\begin{minted}{julia}
using DataFrames

# Create a DataFrame from columns
df = DataFrame(
    name   = ["Alice", "Bob", "Charlie", "Diana", "Eve"],
    age    = [28, 35, 42, 31, 27],
    salary = [55000, 72000, 88000, 63000, 51000],
    dept   = ["Engineering", "Marketing", "Engineering", "HR", "Marketing"]
)

# Basic inspection
size(df)              # (5, 4)
nrow(df), ncol(df)    # (5, 4)
names(df)             # ["name", "age", "salary", "dept"]
describe(df)          # summary statistics
first(df, 3)          # first 3 rows
last(df, 2)           # last 2 rows
eltype.(eachcol(df))  # column types
\end{minted}

% ============================================================
\section{Reading and writing data}

\begin{minted}{julia}
using CSV, DataFrames

# Read CSV files
df = CSV.read("data.csv", DataFrame)

# Read with options
df = CSV.read("data.csv", DataFrame;
    delim=',',
    header=true,
    missingstring="NA",
    types=Dict(:age => Int, :salary => Float64),
    select=[:name, :age, :salary]   # read only these columns
)

# Read directly from a URL
using Downloads
url = "https://raw.githubusercontent.com/owid/covid-19-data/" *
      "master/public/data/owid-covid-data.csv"
Downloads.download(url, "covid.csv")
covid = CSV.read("covid.csv", DataFrame)

# Write to CSV
CSV.write("output.csv", df)

# Read from RDatasets (built-in classic datasets)
using RDatasets
iris = dataset("datasets", "iris")
mtcars = dataset("datasets", "mtcars")
\end{minted}

\begin{juliatip}
\texttt{CSV.read} is very fast --- often faster than pandas for large files. For even better performance on huge files, use \texttt{CSV.read("big.csv", DataFrame; ntasks=8)} to parallelize parsing across threads.
\end{juliatip}

% ============================================================
\section{Selecting and filtering}

\begin{minted}{julia}
using DataFrames, RDatasets

iris = dataset("datasets", "iris")

# Select columns
select(iris, :SepalLength, :Species)
select(iris, Not(:PetalWidth))
select(iris, Between(:SepalLength, :PetalLength))
select(iris, r"Sepal")          # regex column selection

# Create new columns with select
select(iris, :SepalLength, :SepalWidth,
    [:SepalLength, :SepalWidth] => ByRow((l, w) -> l * w) => :SepalArea
)

# Filter rows
filter(:Species => ==("setosa"), iris)
filter(row -> row.SepalLength > 5.0 && row.Species == "setosa", iris)
subset(iris, :SepalLength => x -> x .> 5.0)

# Combine selection and filtering
iris |>
    x -> filter(:Species => ==("virginica"), x) |>
    x -> select(x, :SepalLength, :PetalLength)
\end{minted}

% ============================================================
\section{Transforming data}

\begin{minted}{julia}
using DataFrames, Statistics

df = DataFrame(
    city = ["Paris", "London", "Berlin", "Madrid", "Rome"],
    pop_million = [2.16, 8.98, 3.64, 3.22, 2.87],
    area_km2 = [105, 1572, 892, 604, 1285]
)

# Add a new column
transform(df, [:pop_million, :area_km2] =>
    ByRow((p, a) -> p * 1e6 / a) => :density)

# In-place transformation
transform!(df, :city => ByRow(uppercase) => :city_upper)

# Multiple transformations
transform(df,
    :pop_million => (x -> x .* 1e6) => :population,
    :area_km2 => (x -> x ./ 2.59) => :area_sq_miles
)

# Replace values
df.city = replace.(df.city, "Paris" => "Paris (France)")

# Handle missing data
df_missing = DataFrame(
    x = [1, 2, missing, 4, missing],
    y = [missing, 2, 3, missing, 5]
)
dropmissing(df_missing)          # drop rows with any missing
dropmissing(df_missing, :x)      # drop rows where x is missing
coalesce.(df_missing.x, 0)      # replace missing with 0
\end{minted}

% ============================================================
\section{Sorting}

\begin{minted}{julia}
using DataFrames, RDatasets

mtcars = dataset("datasets", "mtcars")

# Sort by one column
sort(mtcars, :MPG)                    # ascending
sort(mtcars, :MPG, rev=true)          # descending

# Sort by multiple columns
sort(mtcars, [:Cyl, :MPG])            # Cyl ascending, then MPG ascending
sort(mtcars, [:Cyl, order(:MPG, rev=true)])

# In-place sort
sort!(mtcars, :HP, rev=true)
\end{minted}

% ============================================================
\section{Grouping and split-apply-combine}

\begin{minted}{julia}
using DataFrames, Statistics, RDatasets

iris = dataset("datasets", "iris")

# Group by species
gdf = groupby(iris, :Species)
length(gdf)    # 3 groups

# Apply aggregation
combine(gdf,
    :SepalLength => mean => :mean_sepal_length,
    :SepalLength => std  => :std_sepal_length,
    :PetalLength => mean => :mean_petal_length,
    nrow => :count
)
# 3×5 DataFrame: one row per species

# Multiple grouping columns
mtcars = dataset("datasets", "mtcars")
combine(groupby(mtcars, [:Cyl, :Gear]),
    :MPG => mean => :avg_mpg,
    :HP  => mean => :avg_hp,
    nrow => :count
)

# Transform within groups (keeps all rows)
transform(groupby(iris, :Species),
    :SepalLength => (x -> (x .- mean(x)) ./ std(x)) => :SepalLength_zscore
)

# Filter groups
subset(groupby(mtcars, :Cyl),
    :MPG => x -> x .> median(x)   # keep above-median MPG within each Cyl group
)
\end{minted}

\begin{juliatip}
The split-apply-combine pattern in DataFrames.jl is: (1) \texttt{groupby} to split, (2) \texttt{combine}/\texttt{transform}/\texttt{subset} to apply and combine. This is analogous to pandas \texttt{groupby().agg()} or dplyr \texttt{group\_by() \%>\% summarise()}.
\end{juliatip}

% ============================================================
\section{Joins}

\begin{minted}{julia}
using DataFrames

# Two related tables
employees = DataFrame(
    id = [1, 2, 3, 4, 5],
    name = ["Alice", "Bob", "Charlie", "Diana", "Eve"],
    dept_id = [10, 20, 10, 30, 20]
)

departments = DataFrame(
    dept_id = [10, 20, 30, 40],
    dept_name = ["Engineering", "Marketing", "HR", "Finance"]
)

# Inner join (only matching rows)
innerjoin(employees, departments, on=:dept_id)

# Left join (keep all employees)
leftjoin(employees, departments, on=:dept_id)

# Right join (keep all departments)
rightjoin(employees, departments, on=:dept_id)

# Outer join (keep everything)
outerjoin(employees, departments, on=:dept_id)

# Join on different column names
sales = DataFrame(emp_id=[1,2,3], revenue=[1000, 2000, 1500])
innerjoin(employees, sales, on=:id => :emp_id)

# Anti-join (employees NOT in sales)
antijoin(employees, sales, on=:id => :emp_id)

# Semi-join (employees who ARE in sales, but only employee columns)
semijoin(employees, sales, on=:id => :emp_id)
\end{minted}

% ============================================================
\section{Reshaping: stack and unstack}

\begin{minted}{julia}
using DataFrames

# Wide format
wide = DataFrame(
    country = ["France", "Germany", "Italy"],
    pop_2020 = [67.39, 83.24, 59.55],
    pop_2021 = [67.75, 83.16, 59.24],
    pop_2022 = [67.97, 84.08, 58.86]
)

# Wide → Long (stack / melt)
long = stack(wide, r"pop_", :country;
    variable_name=:year, value_name=:population)

# Clean the year column
transform!(long, :year => ByRow(y -> parse(Int, last(string(y), 4))) => :year)

# Long → Wide (unstack)
unstack(long, :country, :year, :population)
\end{minted}

\begin{exercise}
\begin{enumerate}
  \item Load the iris dataset from RDatasets. Compute the mean and standard deviation of all four measurements, grouped by species.
  \item Load the mtcars dataset. Filter for cars with more than 100 HP and 6 or more cylinders. Sort by MPG descending.
  \item Download the Our World in Data COVID-19 CSV. Filter for a country of your choice, select date and new cases, and compute the 7-day rolling average.
  \item Create two DataFrames: \texttt{students} (id, name, major) and \texttt{grades} (student\_id, course, grade). Perform an inner join and compute the average grade per major.
  \item Using the mtcars dataset, demonstrate the split-apply-combine pattern: group by number of cylinders, compute mean MPG, mean HP, and count.
  \item Create a wide-format DataFrame with monthly sales for 4 products over 12 months. Reshape to long format, then compute total annual sales per product.
\end{enumerate}
\end{exercise}

% ============================================================
\section{Chapter summary}

\begin{itemize}
  \item DataFrames.jl is Julia's tabular data package, similar to pandas and dplyr.
  \item CSV.jl reads and writes CSV files efficiently, with parallel parsing.
  \item \texttt{select}, \texttt{filter}/\texttt{subset}, and \texttt{transform} are the core manipulation verbs.
  \item \texttt{groupby} + \texttt{combine} implements the split-apply-combine pattern.
  \item Six join types are available: inner, left, right, outer, semi, anti.
  \item \texttt{stack} (wide to long) and \texttt{unstack} (long to wide) reshape data.
  \item RDatasets.jl provides hundreds of classic datasets for practice.
\end{itemize}
